The following section gives an introduction to what features of Ansible should be supported, how they should be used, and where opinionated defaults may be required.
The first subsection will give an introduction to what should be supported, giving more detail compared to the requirements.
The second subsection shall then define what explicitly should not be implemented at this time, and should be subject to a second version that may be implemented after this thesis is complete.

\subsection{Required features and goals for the implementation}
The following list displays goals that should be fulfilled by the implementation that can be achieved in this thesis.

\subsubsection{Ansible execution}
Execution of Ansible should be handled by Ansible itself.
The Ansible executable should be run against a prebuilt environment that contains all required files for running, authenticating, etc.
This way, there is no need to parse Ansible files, reducing the tightness of coupling between Ansible and the operator, allowing for easier independent updates of Ansible from the operator.

\subsubsection{Inventory management}
The management of inventory objects has already been mentioned in the requirements chapter. In detail, this operator should provide \texttt{CustomResourceDefinitions} for the Ansible inventory primitives, and have a mechanism for delivering these primitives to the environment Ansible is executed in.
The inventory management system should also support injecting host- and group-specific variables into the objects, which should be injected into the Ansible environment at runtime.

\subsubsection{Authentication and security}
Security of the objects in this operator themselves will make use of the Kubernetes \gls{rbac} infrastructure to ease the implementation of access controls. Cluster administrators installing the operator into the cluster shall be expected to take care of setting access permissions for users of their cluster.

Security of the \gls{ssh} connection shall be ensured by checking the host keys of the target hosts and injecting them into the Ansible environment at runtime.
This means a \gls{tofu} operation should be performed on the hosts targeted when they are created.
Changed host keys should be identified as a possible \gls{mitm} attack and should be prevented.
For this feature, \gls{ssh} already provides sufficient infrastructure.

Authentication against the hosts should be limited to key-based authentication.
The private keys should be injected into the runtime environment and should be mapped to the \texttt{Hosts}.
Each \texttt{Host} may have a different key that needs to be used, so there needs to be a mapping for each \texttt{Host} with a separate private key file.

\subsubsection{Playbook sourcing and dependencies}
Playbooks should both be able to be inline for easy management within Kubernetes, and be able to be pulled from Git repositories so they can be easier version controlled.
Pulling playbooks from Git should be possible via \gls{http}.
Playbooks should also have the ability to have dependencies via a \enquote{requirements.yml} file that should be resolved by the operator.

\subsection{Explicit non-goals for the implementation}
This subsection lists features that explicitly should not be implemented.

\subsubsection{Ansible Vault}
Ansible Vault is a way of injecting confidential variables into an Ansible inventory in the form of encrypted files which allows Ansible to receive confidential data via Git or other public means, decrypting the data at runtime.
While this feature is useful, it shall not be supported by the operator at this time, since the implementation would necessitate additional logic.

\subsubsection{Authentication against Git servers and \gls{ssh} communication}
For the time of this thesis, authentication against private Git repositories shall not be implemented.
Implementation would necessitate additional preparation of the runtime environment, Git credential helpers and injecting additional secrets into the runtime.
Similarly, \gls{ssh} communication would necessitate additional \gls{tofu} procedures for secure communication as well as the implementation of credential injection for \gls{ssh} authentication.
